\newcommand{\titulus}{\nomenFesti{S. Augustini, Episcopi \& Ecclesiæ Doctoris.}
\dies{Die 28. Augusti.}}
\newcommand{\oratio}{\pars{Oratio.}

\noindent Innova, quǽsumus, Dómine, in Ecclésia tua spíritum quo beátum Augustínum, epíscopum, imbuísti, ut, eódem nos repléti, te solum veræ fontem sapiéntiæ sitiámus et supérni amóris quærámus auctórem.

\pars{Pro pace in universo mundo.} \scriptura{Sir. 50, 25; 2 Esdr. 4, 20; \textbf{H416}}

\vspace{-4mm}

\antiphona{II D}{temporalia/ant-dapacemdomine.gtex}

\vfill

\noindent Deus, a quo sancta desidéria, recta consília et iusta sunt ópera: da servis tuis illam, quam mundus dare non potest, pacem; ut et corda nostra mandátis tuis dédita, et hóstium subláta formídine, témpora sint tua protectióne tranquílla.

\noindent Per Dóminum nostrum Iesum Christum, Fílium tuum, qui tecum vivit et regnat in unitáte Spíritus Sancti, Deus, per ómnia sǽcula sæculórum.

\noindent \Rbardot{} Amen.}
\newcommand{\invitatorium}{\pars{Invitatorium.}

\vspace{-4mm}

\antiphona{IV}{temporalia/inv-fontemsapientiae.gtex}}
\newcommand{\hymnusmatutinum}{\pars{Hymnus}

\cuminitiali{I}{temporalia/hym-MagnePater.gtex}}
\newcommand{\matversus}{\noindent \Vbardot{} Audies de ore meo verbum.

\noindent \Rbardot{} Et annuntiábis eis ex me.}
\newcommand{\lectioi}{\pars{Lectio I.} \scriptura{Ier. 4, 5-8. 13-19}

\noindent De libro Ieremíæ prophétæ.

\noindent Hæc dicit Dóminus:

\noindent «Annuntiáte in Iuda et in Ierúsalem audítum fácite; et loquímini et cánite tuba in terra, clamáte fórtiter et dícite: “Congregámini, et ingrediámur civitátes munítas”. Leváte signum in Sion, fúgite, nolíte stare, quia malum ego addúco ab aquilóne et contritiónem magnam. Ascéndit leo de cubíli suo et prædo géntium se levávit; egréssus est de loco suo, ut ponat terram tuam in solitúdinem: civitátes tuæ vastabúntur, remanéntes absque habitatóre. Super hoc accíngite vos cilíciis, plángite et ululáte, quia non est avérsa ira furóris Dómini a nobis.

\noindent Ecce quasi nubes ascéndet, et quasi tempéstas currus eius; velocióres áquilis equi illíus. Væ nobis, quóniam vastáti sumus! Lava a malítia cor tuum, Ierúsalem, ut salva fias; úsquequo morabúntur in te cogitatiónes iníquæ? Vox enim annuntiántis a Dan et notam faciéntis calamitátem de monte Ephraim. Nuntiáte géntibus. Ecce adsunt! Audítum fácite hoc super Ierúsalem: «Custódes venérunt de terra longínqua et dedérunt super civitátes Iudæ vocem suam; quasi custódes agrórum facti sunt super eam in gyro, quia advérsus me cóntumax erat», dicit Dóminus.

\noindent Via tua et ópera tua fecérunt hæc tibi; ista malítia tua, quia amára, quia tétigit cor tuum. Víscera mea, víscera mea! Dóleo. Paríetes cordis mei! Turbátur in me cor meum: non tacébo, quóniam vocem búcinæ audívit ánima mea, clamórem prœ́lii.}
\newcommand{\responsoriumi}{\pars{Responsorium 1.} \scriptura{\Rbardot{} Ier. 4, 26.24.27 \Vbardot{} Ps. 8, 2; \textbf{H418}}

\vspace{-5mm}

\responsorium{II}{temporalia/resp-afaciefuroristui-CROCHU.gtex}{}}
\newcommand{\lectioii}{\pars{Lectio II.} \scriptura{Ier. 4, 20-28}

\noindent Contrítio super contritiónem vocáta est, quóniam vastáta est omnis terra, repénte vastáta sunt tabernácula mea, súbito tentória mea. Usquequo vidébo vexíllum, áudiam vocem búcinæ? «Quia stultus pópulus meus: me non cognovérunt; fílii insipiéntes sunt et vecórdes: sapiéntes sunt, ut fáciant mala, bene autem fácere nésciunt».

\noindent Aspéxi terram, et ecce vácua erat et desérta; et cælos, et non erat lux in eis.

\noindent Aspéxi montes, et ecce movebántur, et omnes colles conturbáti sunt.

\noindent Aspéxi, et ecce non erat homo, et omne volátile cæli recésserat.

\noindent Aspéxi, et ecce hortus desértus, et omnes urbes eius destrúctæ sunt a fácie Dómini et a fácie iræ furóris eius.

\noindent Hæc enim dicit Dóminus: «Desérta erit omnis terra, sed tamen consummatiónem non fáciam. Super hoc lugébit terra, et mærébunt cæli désuper, eo quod locútus sum, státui et non pǽnitet me, nec avértar ab eo».}
\newcommand{\responsoriumii}{\pars{Responsorium 2.} \scriptura{\Rbardot{} Mic. 6, 3; \textbf{H161}}

\vspace{-5mm}

\responsorium{VII}{temporalia/resp-populemeusquidfeci-CROCHU.gtex}{}}
\newcommand{\lectioiii}{\pars{Lectio III.} \scriptura{Lib. 7, 10. 18; 10, 27: CSEL 33, 157-163. 255}

\noindent Ex Confessiónum libris sancti Augustíni epíscopi.

\noindent Admónitus redíre ad memetípsum, intrávi in íntima mea duce te et pótui, \emph{quóniam factus es adiútor meus.} Intrávi et vidi qualicúmque óculo ánimæ meæ supra eúndem óculum ánimæ meæ, supra mentem meam lucem incommutábilem, non hanc vulgárem et conspícuam omni carni nec quasi ex eódem génere grándior erat, tamquam si ista multo multóque clárius clarésceret totúmque occupáret magnitúdine. Non hoc illa erat, sed áliud, áliud valde ab istis ómnibus. Nec ita erat supra mentem meam, sicut óleum super aquam nec sicut cælum super terram, sed supérior, quia ipsa fecit me, et ego inférior, quia factus ab ea. Qui novit veritátem, novit eam.

\noindent O ætérna véritas et vera cáritas et cara ætérnitas! Tu es Deus meus, tibi suspíro die ac nocte. Et cum te primum cognóvi, tu assumpsísti me, ut vidérem esse, quod vidérem, et nondum me esse, qui vidérem. Et reverberásti infirmitátem aspéctus mei rádians in me veheménter, et contrémui amóre et horróre; et invéni longe me esse a te in regióne dissimilitúdinis, tamquam audírem vocem tuam de excélso: «Cibus sum grándium: cresce et manducábis me. Nec tu me in te mutábis sicut cibum carnis tuæ, sed tu mutáberis in me».

\noindent Et quærébam viam comparándi róboris; quod esset idóneum ad fruéndum te, nec inveniébam, donec amplécterer \emph{mediatórem Dei et hóminum, hóminem Christum Iesum, qui est super ómnia Deus benedíctus in sǽcula,} vocántem et dicéntem: \emph{Ego sum via et véritas et vita,} et cibum, cui capiéndo inválidus eram, miscéntem carni, quóniam \emph{Verbum caro factum est,} ut infántiæ nostræ lactésceret sapiéntia tua, per quam creásti ómnia.

\noindent Sero te amávi, pulchritúdo tam antíqua et tam nova, sero te amávi! Et ecce intus eras et ego foris et ibi te quærébam et in ista formósa, quæ fecísti, defórmis irruébam. Mecum eras, et tecum non eram. Ea me tenébant longe a te, quæ si in te non essent, non essent. Vocásti et clamásti et rupísti surditátem meam, coruscásti, splenduísti et fugásti cæcitátem meam, fragrásti, et duxi spíritum et anhélo tibi, gustávi et esúrio et sítio, tetigísti me, et exársi in pacem tuam.}
\newcommand{\responsoriumiii}{\pars{Responsorium 3.} \scriptura{\Rbardot{} Sir. 15, 3 \Vbardot{} Ps. 131, 18; \textbf{H62}}

\vspace{-5mm}

\responsorium{VII}{temporalia/resp-cibavitillum-CROCHU-cumdox.gtex}{}}
\newcommand{\hymnuslaudes}{\pars{Hymnus.}

\cuminitiali{I}{temporalia/hym-FulgetInCaelisAugustino.gtex}}
\newcommand{\lectiobrevis}{\pars{Lectio Brevis.} \scriptura{Sap. 7, 13-14}

\noindent Sapiéntiam sine fictióne dídici et sine invídia commúnico; divítias illíus non abscóndo. Infinítus enim thesáurus est homínibus; quem qui acquisiérunt, ad amicítiam in Deum se paravérunt propter disciplínæ dona commendáti.}
\newcommand{\responsoriumbreve}{\pars{Responsorium breve.} \scriptura{Sir. 45, 9}

\antiphona{VI}{temporalia/resp-amaviteum.gtex}}
\newcommand{\benedictus}{\pars{Canticum Zachariæ.}

\vspace{-4mm}

\antiphona{VIII G}{temporalia/ant-tuesmagistermeus.gtex}

\vspace{-2mm}

\scriptura{Lc. 1, 68-79}

\vspace{-2mm}

\cantusSineNeumas
\initiumpsalmi{temporalia/benedictus-initium-viii-G-auto.gtex}

%\vspace{-1.5mm}

\input{temporalia/benedictus-viii-G.tex} \Abardot{}}
\newcommand{\preces}{\noindent Christo, bono pastóri,~\gredagger{} qui pro suis óvibus ánimam pósuit,~\grestar{} laudes grati exsolvámus et supplicémus, dicéntes:

\Rbardot{} Pasce pópulum tuum, Dómine.

\noindent Christe, qui in sanctis pastóribus misericórdiam et dilectiónem tuam dignátus es osténdere,~\grestar{} numquam désinas per eos nobíscum misericórditer ágere.

\Rbardot{} Pasce pópulum tuum, Dómine.

\noindent Qui múnere pastóris animárum fungi per tuos vicários pergis,~\grestar{} ne destíteris nos ipse per rectóres nostros dirígere.

\Rbardot{} Pasce pópulum tuum, Dómine.

\noindent Qui in sanctis tuis, populórum dúcibus, córporum animarúmque médicus exstitísti,~\grestar{} numquam cesses ministérium in nos vitæ et sanctitátis perágere.

\Rbardot{} Pasce pópulum tuum, Dómine.

\noindent Qui, prudéntia et caritáte sanctórum, tuum gregem erudísti,~\grestar{} nos in sanctitáte iúgiter per pastóres nostros ædífica.

\Rbardot{} Pasce pópulum tuum, Dómine.}
\newcommand{\benedicamuslaudes}{\cuminitiali{}{temporalia/benedicamus-memoria-laudes.gtex}}
\include{hebdomadaxxi}
\include{feriavi}
